\chapter{Sprint 56 --- Batch C1--C10: đóng 32 issue \texttt{Ready for Dev} của tab \texttt{5. Issue List} (2026-08-15 $\rightarrow$ 2026-08-23)}

\section{Bối cảnh}

Ngày 14/08/2026 tab \texttt{5. Issue List} tồn \textbf{32 issue} ở trạng thái
\texttt{Ready for Dev} (STT 83--114), toàn bộ \texttt{Owner = Dev},
\texttt{Need BA Confirm = No}. Phiên 14/08 \textbf{không viết một dòng code nào}: đọc hết 32
issue, đối chiếu với source thật, rồi chia thành \textbf{10 cụm theo gốc rễ chung} chứ không
theo thứ tự trong sheet --- bảng cụm nằm ở \texttt{docs/next-session-clusters-C.md} và §13 của
compact summary.

Chủ dự án chốt ba luật cho cả batch:

\begin{enumerate}
\item \textbf{Mỗi phiên đúng một cụm}, một lần \texttt{-u}, một lần deploy UAT. Không gộp cụm
  để ``cho nhanh''.
\item \textbf{Không được chạy \texttt{/wujia-end-sprint} cho tới khi đủ 10/10 cụm.} Chốt sprint
  sớm thì chapter và commit lệch nhau, phiên sau nạp ngữ cảnh sai.
\item \textbf{Clean code}: sửa ở đúng chỗ gốc, tách helper dùng chung thay vì chép; không thêm
  field/file khi cái sẵn có đủ; comment chỉ giải thích \emph{tại sao} ở chỗ khó đoán.
\end{enumerate}

Thứ tự chạy \textbf{không phải} C1$\rightarrow$C10 tuần tự mà là
C1 $\rightarrow$ C2 $\rightarrow$ C3 $\rightarrow$ C4 $\rightarrow$ C5 $\rightarrow$ C9
$\rightarrow$ C10 $\rightarrow$ C6 $\rightarrow$ C7 $\rightarrow$ C8, vì có ràng buộc phụ
thuộc: C1 sửa gốc dữ liệu mà C2/C3/C5 cần để retest; C6 (trạng thái tương tác của component
chung) phải xong trước C7 (Home dùng lại các trạng thái đó); C8 chạm heading của mọi trang nên
để cuối, chạy sau khi C7 đã sắp lại Home thì khỏi làm hai lần.

\section{Kết quả từng cụm}

\begin{itemize}
\item \textbf{C1 --- Franchise backend} (3 issue, \texttt{b623b70}, 15/08). Gốc rễ: cửa hàng
  không suy được từ đối tác, không được kiểm tra ở backend, và không kế thừa khi tạo giấy báo
  có. Một helper dùng chung partner $\rightarrow$ cửa hàng (một truy vấn có chỉ mục, không tìm
  trong vòng lặp), chặn ở \texttt{action\_confirm} / \texttt{button\_validate} /
  \texttt{action\_post} nên import và API cũng không lọt.

\item \textbf{C2 --- Công nợ, luồng dữ liệu} (4 issue, \texttt{c656b95}, 15/08). 38/38 ô đo
  đạt, 113 test xanh.

\item \textbf{C3 --- Công nợ, giao diện} (5 issue, \texttt{0398fa3}, 15/08). 115/115 ô đo đạt
  ở 6 chiều rộng, 38 test $+$ 116 hồi quy xanh.

\item \textbf{C4 --- Kiến thức} (4 issue, \texttt{a292547}, 15/08). 60/60 ô đo, 14 test mới
  $+$ 87 hồi quy.

\item \textbf{C5 --- Giao hàng} (4 issue, \texttt{94756ff}, 15/08). 20/20 ô đo, 9 test mới
  $+$ 75 hồi quy.

\item \textbf{C9 --- Lịch sử đặt hàng} (3 issue, \texttt{9a8784b}, 15/08). 26/26 ô đo ở
  360/391/1920. Bẫy: CSS gốc nằm ở \texttt{wujia\_portal\_layout} nên lệnh \texttt{-u} phải
  kèm module đó, thiếu là deploy xong không thấy đổi gì.

\item \textbf{C10 --- Đăng ký thi} (2 issue, \texttt{f99321b}, 18/08). 33/33 ô đo, 20 test mới
  $+$ 82 hồi quy. Đã deploy UAT và đo lại trên UAT 30/30 --- trừ nhánh ``ca thi thắng khoá'' vì
  UAT hết ca mở.

\item \textbf{C6 --- Trạng thái tương tác} (1 issue, \texttt{c2592b3}, 19/08). Xem §3.

\item \textbf{C7 --- Trang chủ} (5 issue, \texttt{b33b46f}, 19/08). Năm issue \textbf{một gốc
  rễ}: Home dựng dồn qua nhiều sprint nên mỗi khối một khung chứa, một dấu hiệu bấm được, một
  giới hạn số dòng riêng. Gom về một hằng \texttt{HOME\_PREVIEW\_LIMIT} và một bộ
  \texttt{.wujia-mdash-card/-row}; 11 mũi tên cấp dòng $\rightarrow$ 0; chiều cao Home
  2752 $\rightarrow$ 2634 điểm ảnh ở khổ 360 với cùng dữ liệu.

\item \textbf{C8 --- Tiêu đề khối danh sách} (1 issue, hai phiên: \texttt{59299f1} 22/08 và
  \texttt{02f77ca} 23/08). Xem §4.
\end{itemize}

\section{C6 --- gốc rễ nằm ở thứ tự nạp CSS, không nằm ở CSS của mình}

Issue báo: mọi liên kết trong portal bị gạch chân khi rê chuột, và nhiều chỗ mất viền chỉ dấu
bàn phím. Đọc code custom thì \emph{không có} rule nào gạch chân. Gốc rễ thật là \textbf{một
rule \texttt{a:hover\{text-decoration:underline\}} của bundle \texttt{web.assets\_frontend}
được nạp SAU toàn bộ CSS custom}.

Cách xử lý: thay vì rải \texttt{!important} khắp nơi, gom toàn bộ trạng thái tương tác vào một
tầng mới \texttt{\_interaction.css} nạp cuối cùng và thắng bằng độ đặc hiệu của bộ chọn.

Điểm đáng ghi nhất của cụm này là \textbf{local không tái hiện được lỗi}: UAT có cài
\texttt{website\_sale} nên bundle khác hẳn máy dev. Bằng chứng phải lấy bằng cách nhúng thẳng
CSS ứng viên vào chính UAT rồi đo trước/sau: gạch chân 18$\rightarrow$0 ở khổ 1920 và
34$\rightarrow$0 ở khổ 391, thiếu viền chỉ dấu 46$\rightarrow$0 và 43$\rightarrow$0, 0 dịch
layout; đi phím Tab thật 7/124 $\rightarrow$ 124/124.

\section{C8 --- một khuôn duy nhất cho tiêu đề khối danh sách}

Đây là cụm dài nhất, phải tách làm hai phiên.

\subsection{Kiểm kê trước, code sau}

Portal có \textbf{183 tiêu đề thô}. Phân loại \textbf{theo tổ tiên trong cây DOM} chứ không
theo tên class: nếu tổ tiên gần nhất có nền $+$ bo góc $+$ đệm thì đó là tiêu đề của thẻ, thuộc
component khác, ngoài phạm vi. Kết quả: 89 tiêu đề thẻ $+$ 3 tiêu đề trang $+$ 24 loại khác
nằm ngoài phạm vi, còn lại là các chỗ thật sự phải đổi
(\texttt{docs/c8-heading-inventory.md}).

\subsection{Hai điều Odoo 19 không cho làm}

\begin{itemize}
\item \textbf{QWeb Odoo 19 không có directive đặt tên thẻ động.} Cặp
  \texttt{t-tag-open/t-tag-close} là chi tiết kỹ thuật bên trong trình biên dịch, sinh ra từ
  thẻ tĩnh (\texttt{ir\_qweb.py:1705}), không dùng được từ template. Vì vậy component phải rẽ
  ba nhánh \texttt{t-if/t-elif/t-else} cho ba cấp tiêu đề --- xấu nhưng là cách duy nhất.
\item \textbf{Comment XML cấm hai gạch nối liền nhau.} Viết tên một modifier BEM có
  \texttt{-{}-} vào comment là chết build với \texttt{XMLSyntaxError}. Phải diễn đạt bằng lời.
\end{itemize}

\subsection{Ép luật ``tối đa một thứ bên phải'' bằng chính component}

Có màn cần bên phải là con số đếm, có màn cần là liên kết ``Xem tất cả'', và ở trang công nợ
thì \emph{tuỳ dữ liệu} mà là cái nào. Thay vì để mỗi màn tự nhớ luật, component chọn theo thứ
tự ưu tiên \texttt{action > control > meta}: màn cứ truyền cả hai, component tự bỏ bớt.

\begin{verbatim}
<t t-set="_sh_slot"
   t-value="'action' if sh_action_url else
            ('control' if sh_control else ('meta' if sh_meta else 'none'))"/>
\end{verbatim}

\subsection{Chống tái phát}

Sau khi đổi hết, các class tiêu đề rời cũ bị xoá hẳn khỏi mã nguồn, kèm một test quét
\texttt{ir.ui.view.arch\_db} của toàn bộ giao diện và bắt lỗi nếu có class nào sống lại. Ràng
buộc thứ tự: \textbf{chỉ được xoá CSS sau khi đã đổi xong template dùng nó} --- xoá trước là
vỡ giao diện im lặng.

\subsection{Số đo}

C8a 46/46 ô, C8b 62/62 ô trên bản sao cơ sở dữ liệu; 64 test; lưới hồi quy 17 tuyến $\times$ 2
khổ màn hình \textbf{286/286}; đi phím Tab \textbf{447/447} điểm dừng còn viền chỉ dấu. Đo lại
sau khi deploy UAT: 58/60 (hai ô không đo được vì UAT không có hoá đơn đến hạn nên trang thanh
toán chuyển hướng đúng như thiết kế) $+$ 46/46 $+$ 286/286 $+$ 448/448.

\section{Quy trình đo áp cho cả 10 cụm}

Mỗi cụm đều: bản sao cơ sở dữ liệu \textbf{cô lập trên cổng riêng} (không đụng cơ sở dữ liệu
dev chính), một bảng đối chiếu \texttt{Yêu cầu | Đo được | Đạt/Không} phải $\geq$ 90\%, hồi quy
tối thiểu 3 trang $\times$ 2 khổ màn hình, rồi mới ghi ledger và đồng bộ sheet. \textbf{Dev
không tự đóng \texttt{Done}}, chỉ đẩy tới \texttt{Ready for Retest}.

Hai lần trong batch này bảng đo ra đỏ mà \textbf{lỗi nằm ở bộ đo, không nằm ở code}:

\begin{enumerate}
\item Trang công nợ ra rỗng vì bản sao chỉ có đúng một hoá đơn nháp không gắn cửa hàng ---
  mồi dữ liệu xong là đủ điểm.
\item Bốn ô đỏ ở trang hỗ trợ vì phiếu đang mở thuộc cửa hàng khác với cửa hàng đang chọn, nên
  bị đẩy về danh sách \emph{đúng như phân quyền thiết kế}. Đây chính là ``4 ô đỏ có sẵn'' mà
  hai sprint trước từng ghi nhận mà chưa truy ra gốc.
\end{enumerate}

Bài học lặp lại: \textbf{bảng đo ra đỏ thì nghi bộ đo và dữ liệu trước, đừng vội sửa code}.

\section{Gì đã làm (mức nghiệp vụ)}

Batch này dọn 32 lỗi giao diện và dữ liệu mà bên kiểm thử ghi nhận trong đợt rà UAT đầu tháng
8. Người dùng portal --- tức chủ và nhân viên của khoảng 1500 cửa hàng nhượng quyền --- sẽ thấy
portal ``đều tay'' hơn: tiêu đề các khối danh sách cùng một cỡ trên cùng một loại màn hình,
liên kết không còn bị gạch chân loạn khi rê chuột, đi bằng bàn phím luôn nhìn thấy mình đang ở
đâu, trang chủ ngắn lại và các khối cùng một kiểu bấm. Phía kế toán và vận hành thì được lợi ở
phần dữ liệu: đơn hàng, phiếu kho và hoá đơn nay bắt buộc phải gắn đúng cửa hàng, kể cả khi
tạo bằng import hay API, và giấy báo có tự kế thừa cửa hàng của hoá đơn gốc thay vì để trống.

\section{Lý do và đánh đổi}

\begin{itemize}
\item \textbf{Độ đậm tiêu đề giữ 700 thay vì 800.} Bảng thành phần mới ghi 800 nhưng một yêu
  cầu trước đó đang ép cứng toàn portal về 700. Dev \textbf{không tự chọn}: giữ 700 cho khớp
  cái đang chạy, ghi câu hỏi vào phần LIMIT chờ BA chốt.
\item \textbf{Nhịp dọc trang công nợ đổi từ 24/12 sang 16/8.} Con số 24/12 là luật riêng của
  màn đó, từng được BA nghiệm thu ở cụm C3. Chủ dự án chốt bỏ luật riêng để mọi màn cùng một
  cách tính --- ưu tiên đồng bộ hơn giữ nguyên con số cũ. Đã ghi rõ cho BA biết khi retest.
\item \textbf{Khối ``Thông tin chuyển khoản'' đi qua khuôn chung nhưng giữ dáng nhãn nhỏ.} Ép
  đúng cỡ chữ của khuôn thì tràn thẻ cao cố định 150 điểm ảnh. Chọn phương án dung hoà, và giữ
  phần CSS riêng đó \emph{trong} tệp của module công nợ chứ không đẩy vào tệp dùng chung.
\end{itemize}

\section{Còn treo lại cho BA}

Ba câu hỏi cố ý \textbf{không đoán} mà để lại kèm \texttt{Need Clarification}: chọn 800 hay
700; tiêu đề nằm trong khung có tính là tiêu đề khối không (câu này đang chặn đúng một chỗ ở
màn Khảo sát, Dev cố ý chưa đổi); và dòng đếm trên tiêu đề trang có áp cùng quy tắc viết đầy
đủ không. Module Phản hồi khắc phục đã bị gỡ khỏi hệ thống nên \textbf{cố ý} nằm ngoài batch.
